\chapter{Composition}

\label{ch:composition}

This chapter answers how graphs become each other's parts. There are two inclusion mechanisms and
they share one property: both decide everything at \emph{load} time. A \texttt{call} node is
replaced by the callee's nodes; an \texttt{[[attach]]} table weaves a fragment into a host at a
named surface. Nothing is resolved during a run, nothing crosses a process boundary, and there is
no second implementation of \enquote{include}. That single decision is what makes the rest of the
chapter necessary, because it also means the rule that decides where a call returns is a rule about
graph shape rather than a rule about invocation.

The chapter spends most of its length on one such rule, because it is the sharpest edge in the
system. A call returns through every node that declares an \texttt{ok} port and has no outgoing
\texttt{ok} edge --- not through the node the callee names in \texttt{exit}. A \texttt{refuse} node
declares no ports at all, so it is not an exit under that rule, and a callee that ends in a refusal
used to dead-end inside the splice: no edge fired, the run stopped, and the caller's queue row
stayed \texttt{Running} instead of being recorded as parked. Row 364 is the fix, and it is a wiring
rule rather than a second mechanism.

\section{Two inclusions, one shape}

\begin{center}
\footnotesize
\begin{tabular}{@{}p{0.18\textwidth}p{0.36\textwidth}p{0.38\textwidth}@{}}
	\toprule
	 & \textbf{\texttt{call} node} & \textbf{\texttt{[[attach]]} fragment} \\
	\midrule
	What is included & another graph, by name & a fragment file under \texttt{graphs/fragments/} \\
	Naming & \texttt{<call>\_\_<node>} & \texttt{<attach>.<node>}, or \texttt{<attach>@<target>.<node>}
	when the attach spans several targets \\
	Entry & a synthesised head \texttt{<call>\_\_in} & the fragment's own entry, chained \\
	Pin & \texttt{sha256} on the call node & \texttt{pin} on the attach \\
	May contain & calls, nested to depth 4 & no call, no attach, no timer, no store, no budget \\
	Ordering & declaration order of the calls in the file & \texttt{order}, then attach id \\
	Where it runs & wherever its nodes land in the host & a named surface: run start, node before,
	node after, run end \\
	\bottomrule
\end{tabular}
\end{center}

The two are deliberately not interchangeable. A \texttt{call} is the right tool for \enquote{this
graph is a step in a larger process}; an \texttt{attach} is the right tool for \enquote{every run of
this graph should do this first}, or \enquote{every node of this kind should be followed by this
check}. The restriction that a fragment may not call and may not attach is what keeps the weave a
single pass: a fragment is work, never a scheduler and never a composition of its own.

\section{The load, in order}

Every caller that runs or checks a graph goes through one function,
\texttt{load\_host\_with\_graphs}. Its order is fixed, and the order explains which error a bad file
produces.

\begin{center}
\resizebox{\linewidth}{!}{%
\begin{tikzpicture}[node distance=9mm and 12mm]
	\node[b] (src) {graph text};
	\node[b, right=of src] (expand) {expand calls\\\scriptsize depth 4, pin, caps};
	\node[b, right=of expand] (weave) {weave attaches\\\scriptsize pin, surfaces, order};
	\node[b, right=of weave] (load) {parse to types\\\scriptsize grant ceiling};
	\node[b, below=of weave] (validate) {validate\\\scriptsize \texttt{V0}--\texttt{V11}};
	\node[b, left=of validate] (run) {a run or a check};

	\draw[f] (src) -- (expand);
	\draw[f] (expand) -- (weave);
	\draw[f] (weave) -- (load);
	\draw[f] (load) -- (validate);
	\draw[f] (validate) -- (run);
\end{tikzpicture}}
\captionof{figure}{The load pipeline; calls are expanded on the \emph{table} before anything is
typed, which is why an expanded node's id is a string operation rather than a data structure one,
and why \texttt{pii graph show} can print a composed graph at all.}
\end{center}

\begin{enumerate}
	\item \textbf{Expand calls.} The host text is parsed as a TOML table, and each \texttt{kind =
	"call"} node is replaced in place. This happens first so that a call inside a callee is expanded
	by the same code, recursively, with a chain of graph names to detect cycles.
	\item \textbf{Weave attaches.} If the expanded table still has an \texttt{[[attach]]} array, the
	fragments are read, pin-checked, contract-checked and spliced at their surfaces.
	\item \textbf{Parse to types.} The resulting table is loaded into \texttt{GraphSpec}, and the
	registry's grant ceiling is applied here --- the graph and each node are refused if they request
	more than the host allows.
	\item \textbf{Validate.} The static rules run over the expanded graph: ports exist, capabilities
	are within the graph's, cycles are bounded, joins are satisfiable.
\end{enumerate}

The consequence of the order is that a pin error beats a port error. A graph whose callee has
changed reports \texttt{fragment.pin.mismatch} before anyone learns that the graph also has a typo
in an edge --- which is correct, because until the pin is honest the shape of the graph is unknown.

\begin{lstlisting}[language=bash]
pii graph lint graphs/queue.toml     # expand, weave, validate; prints warnings too
pii graph show graphs/queue.toml     # the graph after expansion, node by node
pii graph explain graphs/sdd.toml    # the composition, one level at a time
\end{lstlisting}

\section{A call, before and after}

The clearest way to see the mechanism is to look at the same file twice. Before the load, the queue
carries a \texttt{call} node and three edges that name it. The node is a placeholder: it has a kind,
a graph name, a pin and a caps list, and no ports of its own.

\begin{center}
\resizebox{\linewidth}{!}{%
\begin{tikzpicture}[node distance=9mm and 12mm]
	\node[b] (select) {\texttt{select}};
	\node[b, right=of select] (call) {\texttt{implement}\\\scriptsize \texttt{kind = "call"}};
	\node[b, right=of call] (record) {\texttt{record}};
	\node[p, below=of call] (pin) {pin: sha256 of \texttt{spec.toml}};

	\draw[f] (select) -- node[p,above=2pt,midway]{\texttt{ok}} (call);
	\draw[f] (call) -- node[p,above=2pt,midway]{\texttt{ok}} (record);
	\draw[bk] (call.south east) to[bend left=18] node[p,below=2pt,midway]{\texttt{unverified}} (record.south west);
	\draw[bk] (call.south west) to[bend right=10] node[p,left=2pt,midway]{\texttt{refused}} (record.south);
	\draw[sf] (pin) -- (call);
\end{tikzpicture}}
\captionof{figure}{Before the load: one node standing in for a graph; the two red edges are the ones
row 364 added --- a host may wire a callee's refusal \emph{by name}, which is the whole of the
refusal mechanism.}
\end{center}

After the load the placeholder is gone. Its head is a synthetic node that counts entries, the
callee's nodes are present under namespaced ids, and the host edges have been re-pointed: edges that
arrived at the call now arrive at the head, and edges that left the call now leave the callee's
declared exit.

\begin{center}
\resizebox{\linewidth}{!}{%
\begin{tikzpicture}[node distance=9mm and 13mm]
	\node[b] (select) {\texttt{select}};
	\node[b, right=of select] (in) {\texttt{implement\_\_in}\\\scriptsize counts entries};
	\node[b, right=of in] (mid) {\texttt{code}, \texttt{verify}, \texttt{docs}\\\scriptsize the
	callee's own chain};
	\node[b, right=of mid] (rev) {\texttt{implement\_\_reverify}};
	\node[b, below=of rev] (done) {\texttt{implement\_\_done}};
	\node[b, below=of done] (record) {\texttt{record}};
	\node[b, right=of done] (ref) {\texttt{implement\_\_unverified}\\\scriptsize \texttt{returns}};
	\node[b, right=of ref] (ref2) {\texttt{implement\_\_refused}\\\scriptsize \texttt{returns}};

	\draw[f] (select) -- (in);
	\draw[f] (in) -- (mid);
	\draw[f] (mid) -- (rev);
	\draw[f] (rev) -- node[p,right=2pt,midway]{\texttt{ok}} (done);
	\draw[f] (done) -- node[p,left=2pt,midway]{\texttt{ok}} (record);
	\draw[bk] (rev.east) -- node[p,above=2pt,midway]{\texttt{fail}} (ref);
	\draw[bk] (mid.south) to[bend right=14] node[p,below=2pt,midway]{an \texttt{err} from the chain} (ref2.north);
	\draw[f] (ref) -- (record);
	\draw[f] (ref2.south) -- (record.east);
\end{tikzpicture}}
\captionof{figure}{After the load: the callee's nodes under the call's name space, the transform
that returns on \texttt{ok}, and the two refusal nodes marked \texttt{returns} so the host's own
edges can carry their verdict to \texttt{record}.}
\end{center}

Five details of the expansion are worth knowing because each of them has produced a misunderstanding
at least once:

\begin{itemize}
	\item \textbf{The name space is a string prefix.} A callee node \texttt{code} becomes
	\texttt{implement\_\_code}. There is no nesting structure and no scope; the id is flat, and that is
	why the load refuses an expansion whose ids collide with the host's (\texttt{graph.call.id}).
	\item \textbf{The head is not decoration.} \texttt{<call>\_\_in} exists so that a call can be
	counted: it carries \texttt{max\_entries}, and the run counts completed entries against it. It also
	gives the call a single place to receive its input.
	\item \textbf{Input is one transform per key.} A transform evaluates one top-level expression, so
	a call's \texttt{[node.in]} table becomes a chain of transforms named
	\texttt{<call>\_\_in\_\_<key>}, each after the head. One node could not hold a whole object of
	expressions.
	\item \textbf{The callee's reads are rewritten.} Inside the callee, \texttt{\$.input.x} refers to
	the call's input, and \texttt{\$.nodes.m} refers to the callee's own copies of its nodes. A callee
	cannot read the host's input by accident.
	\item \textbf{The host can read the result by the call's id.} \texttt{\$.nodes.implement} is
	rewritten to \texttt{\$.nodes.implement\_\_done}, so a host expression does not need to know which
	inner node happens to be the exit.
\end{itemize}

That last rewrite has a boundary rule that is easy to get wrong and is tested both ways: the prefix
must end at a dot or at the end of the string. \texttt{\$.nodes.sdd} is rewritten;
\texttt{\$.nodes.sddx} is a different node and is left alone. Node ids and kinds are never rewritten
--- only values --- and a host that never mentions the call id comes out byte-identical.

\section{How a call passes data}

The previous sections show the \emph{shape} of an inclusion. This one answers the question a reader
asks next: when one graph calls another, how does anything travel across the boundary? There are
exactly three moves, and each has a name in the loaded document.

\subsection{In: the caller writes the callee's entry inputs}

A \texttt{call} node carries an \texttt{in} table. Its keys are the \emph{callee's} entry-node inputs,
and its values are expressions resolved in the \emph{caller's} context. From \texttt{queue}, calling
\texttt{spec}:

\begin{verbatim}
[[node]]
id    = "implement"
kind  = "call"
graph = "spec"
sha256 = "771d098b..."

[node.in]
spec   = "$.nodes.select.item.spec"
base   = "$.input.base"
brief  = "Implement the spec named above. Do not commit."
\end{verbatim}

Two of those come from the host's own context and one is a literal. That is the point: the boundary
crosses \emph{values}, not files and not command lines, and each key is computed separately.

At load, every key becomes its own node --- \texttt{implement\_\_in\_\_spec},
\texttt{implement\_\_in\_\_base}, \texttt{implement\_\_in\_\_brief} --- chained behind a head node
\texttt{implement\_\_in}. So an input can be computed, ordered and journalled like any other value;
nothing is special-cased for composition.

\subsection{Inside: the callee is spliced, and its ids are namespaced}

The callee is not a subprocess and not a fresh run. Its nodes are copied into the host document with
ids of the form \texttt{<call>\_\_<node>}, its entry becomes \texttt{<call>\_\_<entry>}, and its exit
becomes \texttt{<call>\_\_<exit>}. Its own edges and ports are carried over unchanged. Nothing about
the callee's internal names survives into the host's vocabulary, which is why a callee may rename its
nodes without breaking a single caller.

\subsection{Out: the host reads the call's own id}

The terminal value of the callee is reached \emph{by the call's id}. The loader rewrites every
occurrence of \texttt{\$.nodes.<call>} in the host --- including sub-paths --- to
\texttt{\$.nodes.<call>\_\_<exit>}. A host that wrote

\begin{verbatim}
brief = "$.nodes.implement.stdout"
\end{verbatim}

is silently read as \texttt{\$.nodes.implement\_\_done.stdout} after the load. The host never names an
internal node; it names the call, and the alias makes the callee's exit answer it. The rewrite is
implemented in \verb|pii-graph/src/fragment/call_output.rs|; the splice and the \texttt{in} transforms
are in \verb|pii-graph/src/fragment/call.rs|.

\begin{center}
\resizebox{\linewidth}{!}{%
\begin{tikzpicture}[node distance=8mm and 12mm]
	\node[b] (host) {\texttt{queue}};
	\node[b, right=of host] (in) {\texttt{implement\_\_in}\\\scriptsize one node per key};
	\node[s, right=of in] (callee) {\texttt{spec}\\\scriptsize spliced inline};
	\node[b, right=of callee] (exit) {\texttt{implement\_\_done}};
	\node[b, right=of exit] (read) {\texttt{\$.nodes.implement}};

	\draw[f] (host) -- node[p,above=2pt,midway]{\texttt{in} keys} (in);
	\draw[f] (in) -- node[p,above=2pt,midway]{\texttt{entry}} (callee);
	\draw[f] (callee) -- node[p,above=2pt,midway]{\texttt{exit}} (exit);
	\draw[sf] (exit) -- node[p,above=2pt,midway]{alias} (read);
\end{tikzpicture}}
\captionof{figure}{The three moves of a call. Inputs are computed in the caller and enter the callee
as \texttt{<call>\_\_in\_\_<key>}; the callee runs inline under namespaced ids; and the host reads the
result by the call's own name, which the loader rewrites to the callee's exit.}
\end{center}

\subsection{Two rules that travel with the boundary}

\begin{enumerate}
	\item A call returns on \emph{its exit's ports}. A callee \texttt{refuse} is pulled back out
	onto the host edges that declared \texttt{ok}, so a host can branch on \texttt{ok} /
	\texttt{unverified} / \texttt{refused} instead of watching the run die.
	\item A call declares \texttt{caps}, and a call may only \textbf{narrow} the host's ceiling. A
	child cannot give itself more than its caller holds, and the ladder is checked at load.
\end{enumerate}

\subsection{The trap: an absent value is not an error}

Values crossing the boundary are untyped JSON. A path that does not exist does \emph{not} fail: it
resolves to nothing, and the caller passes the absence along. Measured on this codebase, a
documentation graph shipped with

\begin{verbatim}
brief = "$.nodes.inventory.payload.stdout"
\end{verbatim}

while a \texttt{tool} node's value is \texttt{\{stdout, ...\}} with no \texttt{payload} level at all.
The brief was empty on every run; the coder was asked to do work with no instructions; and every test
and every gate passed, because none of them executed the graph. The general shape of the fix is to
check paths and placeholders \emph{at load}, when the graph is read, rather than at spawn, when the
money is spent --- and to write at least one test that runs the graph rather than reading it.


\section{The exit rule}

Now the rule the chapter exists for. When a piece of a graph is spliced into another, the splicer
has to answer one question: which of these nodes hands control back? The answer is computed, not
declared:

\begin{enumerate}
	\item For each node in the piece, ask its factory for the ports it declares with this
	configuration.
	\item The node qualifies if that set contains \texttt{ok}.
	\item The node is \emph{wired} if the piece has any outgoing edge from it on port \texttt{ok}.
	\item The piece's exits are the nodes that qualify and are not wired.
\end{enumerate}

The rule is applied twice in the loader, at two scales. Per piece, it decides which nodes feed the
next piece's entry (or the tail the caller supplied). At the host level, \texttt{run.end} uses the
same predicate to decide which host nodes the end fragment follows: every host node with no outgoing
\texttt{ok} edge becomes a sink, and the end chain is hung off all of them with order edges. A graph
with three natural endings therefore reaches the end chain from three places, which is the point ---
a fragment attached at \texttt{run.end} runs however the graph finishes.

\begin{center}
\resizebox{\linewidth}{!}{%
\begin{tikzpicture}[node distance=9mm and 14mm]
	\node[b] (reverify) {\texttt{reverify}\\\scriptsize judge};
	\node[b, right=of reverify] (done) {\texttt{done}\\\scriptsize transform: \texttt{ok}};
	\node[b, below=of reverify] (unv) {\texttt{unverified}\\\scriptsize \texttt{refuse}: no ports};
	\node[b, right=of unv] (ref) {\texttt{refused}\\\scriptsize \texttt{refuse}: no ports};
	\node[b, below=of ref] (errs) {\texttt{code.err}, \texttt{docs.err}\\\scriptsize the other way in};
	\node[v, right=of done] (exit) {is an exit};

	\draw[f] (reverify) -- node[p,above=2pt,midway]{\texttt{ok}} (done);
	\draw[f] (reverify) -- node[p,right=2pt,midway]{\texttt{fail}} (unv);
	\draw[f] (errs) -- (ref);
	\draw[f] (done) -- (exit);
	\draw[bk] (unv.east) -- node[p,below=2pt,midway]{dead end} (ref.west);
\end{tikzpicture}}
\captionof{figure}{The tail of \texttt{graphs/spec.toml} under the exit rule: \texttt{done}
returns because it declares \texttt{ok} and has no \texttt{ok} successor, while \texttt{unverified}
and \texttt{refused} declare nothing, so a refusal leaves no edge out of the splice.}
\end{center}

Two consequences follow, and both are load-bearing.

\paragraph{A piece with no exits is not an error at load.} It is a piece that never returns. The
loader accepts it because a fragment attached at \texttt{run.end} legitimately ends there, and
because a callee's refusal is a legitimate way for the whole run to fail. The failure only becomes
visible at run time, as a run that stops --- which is exactly the class of failure this system tries
to make impossible.

\paragraph{The declared \texttt{exit} is a different thing.} A graph's \texttt{exit} names its
terminal for a direct run, and the call mechanism uses it to re-point host edges: an edge that left
the call now leaves \texttt{<call>\_\_<exit>}. The wider rule exists so that a fragment with several
natural endings still hands control back. Both facts are true at once, which is why reading only one
of them produces a wrong mental model: the declared exit decides where the \emph{host's} edges land,
and the computed exit set decides where a \emph{piece} hands control to what follows it.

\section{Why a refusal dead-ended}

The refusal shape is worth following to the bottom, because it is the one place where a node's port
set is empty rather than merely unused.

\begin{center}
\footnotesize
\begin{tabular}{@{}p{0.30\textwidth}p{0.62\textwidth}@{}}
	\toprule
	\textbf{Fact} & \textbf{Where it is decided} \\
	\midrule
	A \texttt{refuse} node declares no ports & \texttt{RefuseFactory::ports} returns an empty vector
	unless the configuration carries \texttt{returns} \\
	Its \texttt{run} fails the run by name & \texttt{RefuseNode::run} returns a failure whose text is
	\texttt{run.refused\{hook=..., reason=...\}} \\
	An empty port set cannot contain \texttt{ok} & so the exit computation never lists it \\
	No exit means no edge to the tail & so nothing after the callee runs \\
	\bottomrule
\end{tabular}
\end{center}

On 2026-11-04 that combination left a queue row \texttt{Running} for a night. The callee refused,
the run stopped, and the row's claiming run was never recorded as finished --- so the row blocked
the queue, and an operator looking at the queue saw work in progress that had in fact already been
refused. The symptom and the cause were three files apart: the queue's status in
\texttt{pii-store/src/queue.rs}, the edge in \texttt{graphs/queue.toml}, and the port set in
\texttt{pii-graph/src/node.rs}.

There were two honest ways out, and the spec said so explicitly rather than letting the implementer
choose silently:

\begin{enumerate}
	\item \textbf{Make every verdict return on \texttt{ok}.} Give the refusal nodes an \texttt{ok}
	port --- or make them transforms carrying a verdict value --- so that all three tail nodes are
	exits. The host must then distinguish them by the value that arrived, because all three come back
	on the same port.
	\item \textbf{Collapse the tail into one verdict node} with ports \texttt{ok},
	\texttt{unverified} and \texttt{refused}, and let the host wire each to the outcome it wants. This
	makes the callee's outcome vocabulary explicit but changes the callee's tail and every host that
	calls it.
\end{enumerate}

\section{The fix: a wired refusal returns}

The chosen fix is a third shape, and it is the smallest one that keeps both properties: the callee's
tail does not change, and the host gets a distinct port per refusal. It works by letting the
\emph{host} opt in, edge by edge.

The load walks the host's edges. For each edge whose \texttt{from} is the call's id, it takes the
port name and looks for an added node called \texttt{<call>\_\_<port>} whose kind is
\texttt{refuse}. If it finds one, it does two things: it sets \texttt{returns = "<port>"} on that
node, and it re-points the edge to leave that node on \texttt{ok} instead of leaving the call. The
node's factory now reports one port, the exit computation includes it, and its \texttt{run} returns
a verdict instead of failing the run.

\begin{center}
\resizebox{\linewidth}{!}{%
\begin{tikzpicture}[node distance=9mm and 12mm]
	\node[b] (host) {host edge\\\scriptsize \texttt{from = "implement"}\\\scriptsize
	\texttt{port = "unverified"}};
	\node[v, right=of host] (find) {inner node is\\\scriptsize \texttt{refuse}?};
	\node[b, below=of find] (mark) {set \texttt{returns}\\\scriptsize \texttt{ports = [ok]}};
	\node[b, right=of mark] (edge) {re-point the edge\\\scriptsize \texttt{from = <call>\_\_unverified}};
	\node[b, right=of edge] (value) {a verdict value\\\scriptsize \texttt{verdict, port, reason,
	criterion}};
	\node[b, right=of find] (leave) {left for validation};

	\draw[f] (host) -- (find);
	\draw[f] (find) -- node[p,right=2pt,midway]{yes} (mark);
	\draw[f] (mark) -- (edge);
	\draw[f] (edge) -- (value);
	\draw[sf] (find) -- node[p,above=2pt,midway]{no} (leave);
\end{tikzpicture}}
\captionof{figure}{The opt-in, edge by edge; a refusal no host wires is not marked, still declares no
ports, and still fails the run --- which is the control that proves the mechanism is wired and not
merely present.}
\end{center}

The verdict that comes back is a small object, and every field of it is used by a host:

\begin{itemize}
	\item \texttt{verdict} is the literal \texttt{"refused"};
	\item \texttt{port} is the name the host wired (\texttt{unverified}, \texttt{refused}), so a host
	with one \texttt{record} node can tell them apart;
	\item \texttt{reason} is \texttt{"<port>: <declared reason>"}, extended with the arm's own reason
	when the arriving value carried one;
	\item \texttt{criterion} is the first failing criterion from the arriving value, when there is
	one --- which is what lets a parked row name the criterion that failed.
\end{itemize}

The controls that came with the fix are as important as the fix, because a mechanism that fires on
everything is indistinguishable from one that fires correctly. Three near misses are asserted: a
refusal the host does not wire still fails the run; a port that names no refusal node is left on the
call and points at the exit, where validation names it; and the same callee run directly still
fails, because nothing in the direct path sets \texttt{returns}. The callee's own behaviour is
unchanged --- which is the requirement that makes the fix acceptable at all.

The host side is four lines of TOML, and their brevity is the point: the mechanism is visible where
a reviewer reads the graph.

\begin{lstlisting}
[[edge]]
from = "implement"
port = "unverified"
to   = "record"

[[edge]]
from = "implement"
port = "refused"
to   = "record"
\end{lstlisting}

\section{Pins}

A pin is the sha256 of the included file's bytes, written into the file that includes it. A call
carries \texttt{sha256}; an attach carries \texttt{pin}. The two produce different messages because
they are different code paths --- \texttt{fragment.pin.mismatch} from the call expansion,
\texttt{fragment.pin.drift} from the fragment loader --- but they mean the same thing: the file
changed since someone reviewed this inclusion.

\begin{center}
\resizebox{\linewidth}{!}{%
\begin{tikzpicture}[node distance=9mm and 12mm]
	\node[s] (leaf) {\texttt{spec.toml}};
	\node[s, right=of leaf] (mid) {\texttt{queue.toml}};
	\node[s, right=of mid] (top) {\texttt{night.toml}};
	\node[b, below=of mid] (pin1) {pin of \texttt{spec}\\\scriptsize inside queue's bytes};
	\node[b, below=of top] (pin2) {pin of \texttt{queue}\\\scriptsize inside night's bytes};

	\draw[f] (leaf) -- node[p,above=2pt,midway]{hashed into} (mid);
	\draw[f] (mid) -- node[p,above=2pt,midway]{hashed into} (top);
	\draw[sf] (mid) -- (pin1);
	\draw[sf] (top) -- (pin2);
	\draw[bk] (leaf.south) to[bend right=30] node[p,below=2pt,midway]{an edit moves every pin above
	it} (pin2.west);
\end{tikzpicture}}
\captionof{figure}{Pins are transitive by construction rather than by policy: a callee's text contains
its own callee's hash, so editing a leaf changes the middle file's bytes, which changes the top
file's pin.}
\end{center}

That transitivity is not a feature anyone wrote; it falls out of hashing text that contains hashes.
It has one practical consequence worth stating: re-pinning after a leaf change is a mechanical edit
of every file above it, in the same commit, and the diff shows one hash per level with one sentence
of justification. The rule the repository follows is that a pin moves only with a reason in the
commit, because the pin is the only thing that makes \enquote{the night cannot run a graph nobody
reviewed} true.

Row 373 is the cautionary case, and it is not a pin defect at all. \texttt{graphs/night.toml} had a
node whose id was \texttt{sdd}, whose kind was \texttt{call}, and whose \texttt{graph} was
\texttt{queue}. The pin was honest: it pinned the queue. The defect was that the name described
something the composition did not do. The fix was to create the graph the name promised and re-pin
the node to it. A pin protects against a file changing under you; it cannot protect against a file
that was never the right one.

\section{Fragments}

A fragment is a graph file with a \texttt{[fragment]} table instead of a graph's header. It is read
once, at load, and spliced into the host at one of four surfaces.

\begin{center}
\footnotesize
\begin{tabular}{@{}p{0.18\textwidth}p{0.30\textwidth}p{0.44\textwidth}@{}}
	\toprule
	\textbf{Surface} & \textbf{Where the chain goes} & \textbf{What the splice does} \\
	\midrule
	\texttt{run.start} & in front of the host's entry & the chain's exits feed the old entry; the
	host's entry becomes the chain's head \\
	\texttt{node.before} & between a node's inbound edges and the node & inbound edges are re-pointed
	at the chain; the chain's exits feed the node \\
	\texttt{node.after} & on a node's \texttt{ok} edge & the node's \texttt{ok} edges are lifted and
	re-hung after the chain; the node now points at the chain's head \\
	\texttt{run.end} & after every host node with no \texttt{ok} successor & each sink gets an order
	edge to the chain's head; the chain's exits feed nothing \\
	\bottomrule
\end{tabular}
\end{center}

The header is checked before any of that happens. It must declare \texttt{name},
\texttt{inputs}, \texttt{outputs} and a boolean \texttt{may\_block}; a fragment that omits one is
refused with \texttt{fragment.header.missing}. The body must not contain a top-level key that would
give it autonomy --- \texttt{timer}, \texttt{store}, \texttt{budget}, \texttt{limits},
\texttt{schedule}, \texttt{cron} --- nor a node configuration key that would do the same
(\texttt{budget}, \texttt{store}, \texttt{limits}, \texttt{cron}, \texttt{timer\_ms},
\texttt{in\_ms}). The message is \texttt{fragment.autonomy.forbidden}, and its text states the
principle: a fragment is work, never a scheduler.

The block rule is the one that reads as arbitrary until you see the failure it prevents. A fragment
attached at \texttt{run.start} or \texttt{node.before} runs \emph{before} the thing it guards, so a
fragment that can block there can hold the host at a point where nothing else has run. Only a
fragment whose header says \texttt{may\_block = true} may attach at those two surfaces; anything
else is refused with \texttt{fragment.block.invalid}. A fragment at \texttt{node.after} or
\texttt{run.end} runs after work has happened and has no such rule.

\begin{lstlisting}
[fragment]
name      = "record-run"
inputs    = ["run"]
outputs   = ["ok"]
may_block = false
\end{lstlisting}

Naming follows the same flat logic as a call, with one addition. A fragment attached to a single
target is namespaced \texttt{<attach>.<node>}. An attach with \texttt{on = ["a", "b"]} produces one
copy per target, and each copy is namespaced \texttt{<attach>@<target>.<node>}, because two copies
with the same ids would collide. The attach's order is \texttt{order}, then id, and that ordering is
what decides which of two fragments at the same surface runs first. Each inlined node also carries
two stamps in its configuration: \texttt{woven-from}, of the form
\texttt{<fragment>\#<node>}, and \texttt{hook}, the attach id. Those stamps are how a trace of a
woven graph can name the file a node came from, long after the file has changed.

\section{Bounds: depth, cycles, entries}

Three limits keep composition from becoming recursion.

\begin{itemize}
	\item \textbf{Depth 4.} A chain of nested calls deeper than four is refused with
	\texttt{graph.call.depth}, naming the call and the chain. The limit is a constant in the
	expander, not a configuration key, because the right answer for a graph language this small is
	\enquote{shallow}.
	\item \textbf{No cycles.} A call that would re-enter a graph already in the chain is refused with
	\texttt{graph.call.cycle}, and the message prints the chain in order. This is a load-time check,
	not a run-time one: two graphs that call each other cannot be inlined at all.
	\item \textbf{Eight entries.} A call's head counts completed entries, default eight, and the run
	fails once a call is entered more than \texttt{max\_entries} times with \texttt{call 'X' entered
	more than max\_entries (N) times}. The count is cumulative for the run and is reset by no
	enclosing loop, so an inner call inside a looped outer call reaches its own limit after that many
	entries across all the outer passes. An author who needs more raises the inner
\texttt{max\_entries} explicitly, in the file, where the number is reviewable.
\end{itemize}

The entry counter exists because inlining is not looping. A call is textually one copy of the
callee; if the host's shape can reach it repeatedly, the callee runs repeatedly with the same ids,
and without a counter a graph that loops would accumulate state indefinitely. The counter turns
that into a named failure at a stated number.

\section{Every refusal composition can produce}

The table is the chapter in one page: each row is a rule code, the condition that raises it, and the
thing a reader has to look at to fix it. All of them are load-time refusals --- composition has no
run-time errors of its own.

\begin{longtable}{@{}p{0.30\textwidth}p{0.40\textwidth}p{0.22\textwidth}@{}}
	\toprule
	\textbf{Code} & \textbf{Raised when} & \textbf{Look at} \\
	\midrule
	\endhead
	\texttt{graph.call.key} & a call node carries a key outside \texttt{id}, \texttt{kind},
	\texttt{graph}, \texttt{sha256}, \texttt{caps}, \texttt{in}, \texttt{max\_entries} & the call's
	table \\
	\texttt{graph.call.missing} & the call has no \texttt{graph}, \texttt{sha256} or \texttt{caps}
	& the call node \\
	\texttt{graph.call.name} & the graph name is empty, starts with a dot, or contains a path
	separator & the \texttt{graph} value \\
	\texttt{graph.call.not\_found} & no graph directory holds \texttt{<name>.toml} & the graph
	directories given to the load \\
	\texttt{graph.call.read} & the callee exists but cannot be read, or is not UTF-8 & the file \\
	\texttt{fragment.pin.mismatch} & the callee's sha256 differs from the call's & both hashes in the
	message \\
	\texttt{graph.call.exit} & the callee declares no \texttt{exit}, or it names no node & the
	callee's tail \\
	\texttt{graph.call.widen} & a call grants caps the host does not hold, or a callee node needs caps
	the call does not grant & the caps on both sides \\
	\texttt{graph.call.cycle} & the callee is already in the chain & the chain printed in the message
	\\
	\texttt{graph.call.depth} & the chain would nest past four & which call to flatten \\
	\texttt{graph.call.id} & an expanded id collides with a host node & the id named in the message
	\\
	\texttt{graph.call.in} & \texttt{in} is not a table & the call's input block \\
	\texttt{graph.call.max\_entries} & \texttt{max\_entries} is not an integer of at least one & the
	value \\
	\texttt{graph.call.needs} & the callee's \texttt{needs} is malformed & the callee's header \\
	\texttt{graph.call.undeclared} & the callee reads an input its holder does not declare & the
	\texttt{needs} list and the \texttt{in} map \\
	\texttt{graph.call.leak} & a callee reads a host node directly instead of through \texttt{in} &
	the expression named in the message \\
	\texttt{graph.call.no\_dir} & a call exists but the load was given no graph directory & the
	caller's setup \\
	\texttt{fragment.pin.missing} & an attach carries no \texttt{pin} & the attach table \\
	\texttt{fragment.pin.drift} & the fragment's bytes differ from the attach's pin & both hashes \\
	\texttt{fragment.surface.unknown} & the attach names a surface outside the four & the \texttt{at}
	value \\
	\texttt{fragment.target.unknown} & \texttt{on} names a host node that does not exist & the
	\texttt{on} list \\
	\texttt{fragment.surface.unreachable} & a \texttt{run.end} attach has no sink to follow & the
	host's last nodes \\
	\texttt{fragment.header.missing} & the fragment lacks \texttt{name}, \texttt{inputs},
	\texttt{outputs} or a boolean \texttt{may\_block} & the \texttt{[fragment]} table \\
	\texttt{fragment.block.invalid} & a non-blocking fragment attaches at a pre-surface & its
	\texttt{may\_block} \\
	\texttt{fragment.autonomy.forbidden} & the fragment would own a timer, store, schedule or budget
	& the key named in the message \\
	\texttt{fragment.nesting.forbidden} & the fragment contains an \texttt{attach} or a \texttt{call}
	& the fragment body \\
	\texttt{fragment.path.invalid} & the fragment name is empty or contains a path component & the
	\texttt{fragment} value \\
	\bottomrule
\end{longtable}

Two of these are worth recognising on sight because they read as something else.
\texttt{graph.call.undeclared} is not about capabilities: it means the callee reads
\texttt{\$.input.x} for a key nobody mapped, which is the shape of a callee whose contract grew
without its callers being told. And \texttt{graph.call.leak} is the one that protects composition
from becoming implicit coupling --- a callee reading a host node by name would work until the host
renamed it, and would make the pin meaningless as a description of what can change.

\section{What composition cannot do}

Two limits are not bugs and are worth knowing before designing a graph.

\paragraph{There is no iteration primitive.} Composition is static, so a graph cannot say
\enquote{do this call once per criterion}. Everything that needs repetition is one node looping
internally --- a judge running four commands, a coder working through a list --- which is precisely
why the judge has to know how to run criteria at all. A \texttt{foreach} or \texttt{map} over a
list would be a new mechanism, and the repository's rule is that a second inclusion mechanism is
never built. 
\paragraph{Models are per-node strings.} Composition is fine with that, but each piece can carry a
different --- possibly dead --- model id, and a graph that composes three pieces inherits three
chances to name a model that is not served. The model-tier work exists to make that a resolved
mapping rather than a literal per node, and until a graph uses it, the literal is what runs.

\section{A checklist for a new graph}

The list is short and every item on it corresponds to a refusal in this chapter.

\begin{enumerate}
	\item Does a graph already do this? If yes, call it; do not copy it. A copy is a second
	implementation and the two will drift.
	\item Does your graph need to be composed? If so, declare \texttt{exit}. A graph without one is a
	run, not a part.
	\item Grant caps for what your nodes need, and no more. The ladder will catch a mismatch, loudly,
	with the node and the capability named.
	\item Pin every call by sha256 and every attach by pin, and say in the commit why a pin moved.
	\item One responsibility per graph. A loop is a graph; a process is a graph; a schedule is a
	graph.
	\item If a step needs no model and no process, it is a \texttt{transform}, not a coder.
	\item Every failure path gets a port whose name says what happened. If the callee refuses, wire
	the refusal --- otherwise the refusal is invisible to the host.
\end{enumerate}

The last item is the one this chapter was written for. A refusal that is not wired is not an error
the host can see; it is a run that stops.
